\section{Pilot Study I: Bias-Aware Low-Bit Training}
\label{sec:lowbit_pilot}

Low-bit training amplifies the pathological effects of anisotropy.
Under blockwise quantization, scale is determined by extreme elementwise values.
As shown in Section~\ref{subsec:mean_outliers},
these extremes are predominantly mean-dominated.
This suggests that addressing mean bias at its source may substitute
for heavier spectral corrections.

\subsection{Hypothesis}

We test the following hypothesis:

\begin{quote}
\emph{If mean bias is the dominant driver of outliers,
then source-level mean removal should recover most of the stability
benefits attributed to SVD-based spike suppression methods such as Metis.}
\end{quote}

\subsection{Method}

We propose a simple \emph{mean--residual quantization} method, which we call
\emph{Averis} (Averaging-Induced Residual Splitting).
The method targets a concrete empirical observation: for large activation-like matrices in Transformer
training, extreme values (which dominate blockwise quantization scales) are largely driven by a
mean-dominated component that is coherent across tokens.
Accordingly, we explicitly split out a column-mean component via averaging and quantize it separately
from the remaining residual, using only reduction operations and standard quantization kernels.

\paragraph{Notation and quantization operator.}
Let $\mathcal{Q}_b(\cdot)$ denote a $b$-bit quantization operator (e.g., blockwise FP4) applied to a matrix.
We write $\bar{\mathbf{M}}\triangleq \mathcal{Q}_b(\mathbf{M})$ for the quantized version of any matrix
$\mathbf{M}$.
Throughout, $\mathbf{1}$ denotes an all-ones vector with the appropriate dimension.

\paragraph{Forward pass: activation mean--residual splitting.}
Let $\mathbf{W}\in\mathbb{R}^{m\times n}$ be a weight matrix and let
$\mathcal{X}\in\mathbb{R}^{b\times s\times m}$ be an input activation tensor, where $b$ is batch size,
$s$ is sequence length, and $m$ is hidden dimension.
As usual, we reshape $\mathcal{X}$ into a matrix $\mathbf{X}\in\mathbb{R}^{l\times m}$ with $l=b\cdot s$
and compute the GeMM output
\[
\mathbf{Y}=\mathbf{X}\mathbf{W}\in\mathbb{R}^{l\times n}.
\]

We define the \emph{activation column-mean vector} $\mu_{\mathbf{X}}\in\mathbb{R}^{m}$ and the associated
\emph{mean matrix} $\mathbf{M}_{\mathbf{X}}\in\mathbb{R}^{l\times m}$ by
\[
\mu_{\mathbf{X}} \triangleq \frac{1}{l}\mathbf{1}^{\top}\mathbf{X},
\qquad
\mathbf{M}_{\mathbf{X}} \triangleq \mathbf{1}\,\mu_{\mathbf{X}}.
\]
The residual is
\[
\mathbf{X}_{\mathrm{R}} \triangleq \mathbf{X}-\mathbf{M}_{\mathbf{X}}.
\]
We then quantize the two parts independently:
\[
\bar{\mathbf{M}}_{\mathbf{X}}=\mathcal{Q}_b(\mathbf{M}_{\mathbf{X}}),
\qquad
\bar{\mathbf{X}}_{\mathrm{R}}=\mathcal{Q}_b(\mathbf{X}_{\mathrm{R}}),
\qquad
\bar{\mathbf{W}}=\mathcal{Q}_b(\mathbf{W}).
\]
The quantized forward computation is
\begin{equation}
\label{eq:averis_forward}
\hat{\mathbf{Y}}
\;=\;
(\bar{\mathbf{M}}_{\mathbf{X}}+\bar{\mathbf{X}}_{\mathrm{R}})\,\bar{\mathbf{W}}.
\end{equation}
In practice, one may compute $\mu_{\mathbf{X}}$ first, subtract it from $\mathbf{X}$ on the fly to form
$\mathbf{X}_{\mathrm{R}}$, and never explicitly materialize $\mathbf{M}_{\mathbf{X}}$.

\paragraph{Backward pass: output-gradient mean--residual splitting.}
Let $\mathbf{D}\in\mathbb{R}^{l\times n}$ denote the matrix of loss derivatives with respect to
the GeMM output $\mathbf{Y}$, i.e.,
\[
\mathbf{D}\triangleq \frac{\partial \mathcal{L}}{\partial \mathbf{Y}}
\quad (\text{after reshaping}).
\]
We define the \emph{output-gradient column-mean vector} $\mu_{\mathbf{D}}\in\mathbb{R}^{n}$ and mean matrix
$\mathbf{M}_{\mathbf{D}}\in\mathbb{R}^{l\times n}$ by
\[
\mu_{\mathbf{D}} \triangleq \frac{1}{l}\mathbf{1}^{\top}\mathbf{D},
\qquad
\mathbf{M}_{\mathbf{D}} \triangleq \mathbf{1}\,\mu_{\mathbf{D}},
\qquad
\mathbf{D}_{\mathrm{R}} \triangleq \mathbf{D}-\mathbf{M}_{\mathbf{D}}.
\]
Quantize the components independently:
\[
\bar{\mathbf{M}}_{\mathbf{D}}=\mathcal{Q}_b(\mathbf{M}_{\mathbf{D}}),
\qquad
\bar{\mathbf{D}}_{\mathrm{R}}=\mathcal{Q}_b(\mathbf{D}_{\mathrm{R}}).
\]
Then the key GeMMs in backpropagation are evaluated as
\begin{equation}
\label{eq:averis_backward}
\widehat{\frac{\partial \mathcal{L}}{\partial \mathbf{X}}}
\;=\;
(\bar{\mathbf{M}}_{\mathbf{D}}+\bar{\mathbf{D}}_{\mathrm{R}})\,\bar{\mathbf{W}}^{\top},
\qquad
\widehat{\frac{\partial \mathcal{L}}{\partial \mathbf{W}}}
\;=\;
(\bar{\mathbf{M}}_{\mathbf{X}}+\bar{\mathbf{X}}_{\mathrm{R}})^{\top}\,(\bar{\mathbf{M}}_{\mathbf{D}}+\bar{\mathbf{D}}_{\mathrm{R}}).
\end{equation}

\paragraph{Computational profile.}
Averis requires only computing two mean vectors ($\mu_{\mathbf{X}}$ and $\mu_{\mathbf{D}}$) and applying
two subtractions per step, in addition to the baseline quantized GeMMs.
No spectral decomposition or iterative orthogonalization is involved, and all additional operations are
GPU-friendly reductions and elementwise kernels.

\subsection{Experimental Setup}

\label{experiment:setup}
\textbf{Models and Datasets.}
We conduct experiments on GPT-2 (1.1B).
For pretraining, we use the DCLM dataset and train the model on 100B tokens.
The raw data are segmented to split long documents and concatenate short ones, filtered to remove non-Unicode and non-English content.

\textbf{FP4 Training.}
We evaluate the efficacy of Averis under FP4 quantization against both FP4 and BF16 baselines.
All FP4 training in this work adopts W4A4G4 quantization, where weights, activations, and gradients are represented in the E2M1 NVFP4 format.
Stochastic rounding (SR) is applied by default in all FP4 experiments, as it mitigates quantization bias and is orthogonal to Averis.

\textbf{Baselines.}
We compare the following settings:
\begin{itemize}[leftmargin=*,itemsep=2pt]
\item \textbf{BF16}: full-precision baseline.
\item \textbf{Vanilla FP4}: direct NVFP4 quantization (W4A4G4) without structural splitting.
\item \textbf{Averis FP4}: the proposed mean--residual splitting applied to activations and output gradients, with independent quantization of the mean matrices and residuals.
\end{itemize}

\subsection{Results}

\textbf{Training loss.}
Figure~\ref{figure:loss-averis} shows the training loss curves of the BF16 baseline, vanilla FP4, and FP4 + Averis.
Vanilla FP4 exhibits a clear loss gap relative to BF16, whereas Averis consistently narrows this gap throughout training and
achieves a final loss close to the BF16 baseline.

\begin{figure}[t]
    \centering
    \includegraphics[width=0.92\linewidth]{figures/1Bloss.png}
    \caption{Training loss curves for GPT-2 1.1B under BF16, vanilla FP4, and FP4 + Averis.}
    \label{figure:loss-averis}
\end{figure}

\textbf{Downstream performance.}
Table~\ref{table:downstream-task-averis} reports downstream task results for GPT-2 1.1B.
Direct FP4 quantization degrades performance relative to BF16, while FP4 + Averis recovers most of the gap and, in some
metrics, matches the BF16 baseline.

\begin{table}[h]
\caption{Downstream performance on GPT-2 1.1B across different precision settings, reported with task-specific metrics.}
\vspace{-0.7\baselineskip}
\centering
\small
\resizebox{\textwidth}{!}{
\tiny
\begin{tabular}{llccccccccc}
\toprule
Model & Method & Loss & ARC-C & ARC-E & BoolQ & LAMBADA & PIQA & RACE & HellaSwag & Avg \\
\midrule
GPT-2   & BF16          & 3.09 & \textbf{30.7} & 59.5 & \textbf{60.0} & 35.4 & \textbf{64.9} & 47.3 & 41.1 & 48.4 \\
(1.1B)  & FP4           & 3.22 & 29.8 & 57.6 & 59.8 & 35.5 & 63.6 & 46.8 & 39.5 & 47.5 \\
        & FP4 + Averis  & \textbf{3.01} & 30.1 & \textbf{60.0} & 59.5 & \textbf{36.2} & 64.1 & \textbf{48.5} & \textbf{41.9} & \textbf{48.6} \\
\bottomrule
\end{tabular}
}
\label{table:downstream-task-averis}
\end{table}

\subsection{Interpretation}

These results indicate that
low-bit instability is primarily driven by mean-dominated outliers.
Full spectral suppression is not strictly necessary.
Source-level bias conditioning recovers most benefits at a fraction of cost.
